% !TeX root = main_long.tex
\section{RQ3: integration with external data}
\subsection{Method}
\label{sec:m-rq3}

RQ3 asks which externally flagged alterations have sufficient sample depth and which source observations support the answer.
As Listing~\ref{lst:vcf-example} illustrates, \texttt{LAD} values follow the sample's \texttt{LAA} subset rather than the full \texttt{ALT} list.
The full experiment uses four synthetic sites represented both locally and globally, with one sample, converted in the expanded profile; the listing retains two local records and only the fields and alleles needed for the illustrated joins.
Its graph is combined with a separately authored local snapshot of five invented, flagged alterations, sharing no identifiers.
The graph pattern in Listing~\ref{lst:localallele} joins them on contig, position, reference allele, and alternate allele, then retrieves supporting depth through \term{forAllele}.
This assumes the same reference assembly, unchecked against the VCF reference declaration, and identical as-written alleles without normalization or VRS identification.
The sample identifies a VCF column, not a patient.

\begin{lstlisting}[float=htbp,caption={Graph-pattern excerpt from the integration query's \term{WHERE} clause. Prefix declarations, \term{SELECT}, filtering, and provenance retrieval are omitted.},label={lst:localallele}]
# 1. the flagged alteration, from the external annotation snapshot
?annotation a ex:FlaggedAlteration ; ex:contig ?contig ; ex:position ?position ;
            ex:referenceAllele ?ref ; ex:alternateAllele ?alt .
# 2. the same alteration as written in the VCF file
?record vcfc:chrom ?contig ; vcfc:pos ?position ; vcfc:ref ?ref ;
        vcfc:hasCall ?call ; vcfc:hasAltAllele ?altAllele .
?altAllele vcfc:alleleValue ?alt .
# 3. that sample's LAD value, and the declaration giving its cardinality
?call vcfc:hasSampleCall/vcfc:hasFormatValue ?field .
?field vcfc:declaredBy [ vcfc:fieldId "LAD" ; vcfc:fieldNumber ?number ] ;
       vcfc:hasValueItem ?item .
# 4. the depth for THIS allele, via the sample's own local allele set
?item vcfc:forAllele ?altAllele ; vcfc:itemValue ?depth .
\end{lstlisting}

The complete query in \path{RQ3/inputs/local-allele-evidence.rq} selects \texttt{LAD >= 20}, distinguishing correct allele-depth joins from positional misreading while retaining source context.
Expected rows, authored before execution, are compared across eight answer fields. 
The ninth, the source file, is checked against the converted input path.
Each of the three excluded alterations is checked individually.
The example concerns only \texttt{Number=LR} \texttt{LAD}.
Its annotations have no biological meaning and are held locally for deterministic execution.

\subsection{Results}
\label{sec:integration}

The query returns exactly the two expected alterations: \texttt{chr1:1~G>C} at depth 30 and \texttt{chr1:2~A>T} at depth 25.
Each retains the correct file, record, sample, and FORMAT declaration context.
Table~\ref{tab:local-allele-example} follows these joins for the two-record example.
For \texttt{chr1:1~G>C}, the annotation's four location and allele values identify the record's C allele, and \term{forAllele} leads to depth 30 (Figure~\ref{fig:model}).
For \texttt{G>A}, the same location matches, but no LAD item points to A, so the query excludes it.

\begin{table}[htbp]
\centering
\caption{Joining four flagged alterations to the two-record example. A dash means no LAD value for that allele, not zero depth. All numeric values pass the threshold of 20, but the two interpretations return different alterations.}
\label{tab:local-allele-example}
\small
\begin{tabular*}{\linewidth}{@{\extracolsep{\fill}}lrr@{}}
\toprule
Flagged alteration & Through \term{forAllele} & By ALT position (incorrect) \\
\midrule
\texttt{chr1:1 G>C} & 30 & --- \\
\texttt{chr1:1 G>A} & --- & 30 \\
\texttt{chr1:2 A>T} & 25 & --- \\
\texttt{chr1:2 A>C} & --- & 25 \\
\bottomrule
\end{tabular*}
\end{table}

The full experiment also confirms exclusion of \texttt{chr1:3~C>G}, outside the local subset and below threshold even under positional reading.
The incorrect interpretation still returns two plausible rows with provenance.
This counterexample was defined from the fixture and specification before execution; it illustrates the rule, rather than measuring a converter error.
The result is limited to this synthetic example, with no claim of clinical relevance or correctness across all VCF data.
